% 原书第 311 页（PDF 第 324 页）
\BookPage{311}{324}
\BookHeading{郑板桥对佛教的态度}{郑板桥对佛教的态度}{陈进传*}
\BookSubheading{一、前言}

有清一代，诗书画卓然功成者，颇不乏人，但声明亮噪，众所皆知，应属郑板桥，只要提及「扬州八怪」、「诗书画三绝」、「难得糊涂」，几乎无人不晓，广受欢迎，亟思神往。论其艺术成就，可谓超凡脱俗，瑰然出奇，之所以如此，乃因其作品发自本性，流露真情，且能别开生面，独树风格，因而常为后人争相研究的对象。据党明放《郑板桥年谱·附录》所列参考文献的郑板桥论着，就近百册之多。至于期刊论文与扬州八怪专书，尙不计算，否则为数更伙，可谓独占清代文人之鳌头。

上述郑板桥的专着中，除部分是其个人文集外，就是诗书画艺术的探讨，与生平传记的写作，但令人费解的是，郑板桥的作品，常有佛教的身影。可见郑板桥与佛教的关系相当密切，但这方面的探讨，真是凤毛麟角，难得一见。笔者上国家图书馆期刊文献资讯网，查询台湾期刊论文的郑板桥索引篇数，共得八十篇，却无一篇的题目提及佛教。充其量，在其研究行文中，提到郑板桥喜欢跟方外僧侣交往，以及诗画赠与，至于牵涉佛教则大都轻笔带过。金实秋算是少有的例外，他独具慧眼，广泛查阅资料，特于1999年间，撰述〈郑板桥方外友考略〉、〈郑板桥与佛教禅宗〉两文，获得学界的首肯与称誉，后来又陆续增写数篇大作，集结成册，颇有特色，足堪参考。2

\vfill
\begin{BookNotes}
\BookNote{*}{作者任职于佛光大学文化资产与创意学系。}
\BookNote{1.}{党明放，《郑板桥年谱》（北京，首都师范大学出版社，2007年7月），页480-490。}
\BookNote{2.}{金实秋，《郑板桥与佛教禅宗》（北京，宗教文化出版社，2001年4月），页191。}
\end{BookNotes}

% 原书第 312 页（PDF 第 325 页）
\BookPage{312}{325}
尽管郑板桥与佛教渊源深厚，和上人互动甚佳，且常住寺院，进而善用佛语，但郑板桥还是活出自己，依然故我，虽曾自称「板桥居士」，却未信仰佛教，也没皈依佛陀。类此情形，传统中国文人时或有之，可得多例，然彼此所展现的情境与氛围，间有不同，甚至各异其趣，唯郑板桥甚见殊绝之处，深具讨论意义。

\BookSubheading{二、郑板桥的长期佛缘}

魏晋以还的文人都跟佛教难脱关系，即使是反佛、辟佛亦复如此，因反佛者与方外也有往来，如韩愈等，但他们间的关系应不如郑板桥来得密切。郑板桥从小到老都跟佛教频繁接触，其来有自，渊源深厚。

\BookSubheading{（一）家世佛缘}

中国佛寺有的为行善事，照顾贫苦大众，设有义冢或佛塔，以供埋骨。郑板桥数代贫儒，谋生不易，因而归葬佛寺。其祖父郑湜，字清之，十二世长门儒官，殁于康熙三十七年（1698）七月初四日丑时，葬于刹院寺。祖母蔡氏，江苏扬州府兴化县蔡灵皋室女。父之来，字立庵，十三世长门廪生，生于康熙十二年（1673），殁年未详，亦葬于刹院寺。3

关于此事，郑板桥有所回顾，曰：

\begin{BookQuote}
刹院寺祖坟，是东门一枝大家公共的，我因葬父母无地，遂葬其旁。取得风水之力，成进士，作宦数年无恙。是众人之富贵福泽，我一人夺之也，于心安乎？不安乎？〈范县署中寄舍弟书〉
\end{BookQuote}

郑板桥感念及此，后来看到作同样功德的「自在庵」重修时，特加撰述，以示不忘，文曰：

\begin{BookQuote}
庵始於邑侯张公蔚生，廉明慈惠，念水乡穷民棺骨无葬地，于城北九里平望东偏买地为义冢，凡一十二亩三分。即于是庄建佛殿，招僧为住持；固以奉佛，实以修护穷民之冢也。张公去后，佛舍荒、冢地荡，过者伤之。慧圆上人毅然
\end{BookQuote}

\vfill
\begin{BookNotes}
\BookNote{3.}{党明放，《郑板桥年谱》，页3。}
\end{BookNotes}

% 原书第 313 页（PDF 第 326 页）
\BookPage{313}{326}
以重修为己任，众亦敬其素操，翕然从之。…有庵有僧，耕渔之暇，持一畚一锸以修冢，而枯骨于兹有托矣。佛舍修、枯骨聚，而张公仁民爱物之心，传于千古矣。凡庵有兴有废，而是庵泽及枯骨，深得佛理，当久而弗替也。〈自在庵记

郑板桥有一从祖福国，年少出家，乃得道高僧。《兴化佛教通志》记其盛事，曰：

成传，字福国，兴化人，郑板桥从祖，年十六投泰州如来庵雨文大师薙染，十九诣华山受具，参济生和尚于扬之福缘。一夕闻香板声，豁然桶底脱落，乃曰：「击碎疑团见古人，而今不用更参询。头头尽是西来意，法法令彰净妙身。」济颔之，出住泰州龙珠。谢事后步礼五台，回道经赵州古柏林寺，即赵州答庭前柏树子处。檀信一见有缘，坚挽开堂，逐受请进院，自此芳名霭着。庄亲主、刘大人奉旨重新其院，后南旋为济祖编语录，适晋陵华祝虚席，众绅以师补之。未几，复移福缘报法恩也。晚岁退休于邗上圆通庵，乾隆己已春示寂，寿六十，有语录一卷。4

乾隆十年（1745)，郑板桥在范县任上，福国上人到访，祖侄二人异地相逢，僧民两途，兴奋之余，感触亦多，板桥特赋诗，曰：衲衣何日破，四十有余年；白首仍缝绽，青春已结穿。透凉经夏好，等絮入秋便；故友无如此，相看互有怜。〈破被（为从祖福国上人作）〉

及至福国上人盘桓数月，即将离去，郑板桥依依难舍，又作二诗赠行，曰：不向空山卧寂廖，红尘堆里刹竿招；宰官风雨朝停泊，艇子惊呼夜听潮。眼底浮云真幻化，杖头芒𪨗自招摇；悬知法雨无边际，洗净铅华廿四桥。

范城小县无人到，忽漫袈裟暮叩门；一盏寒镫供佛火，数橡茅茨即山村。支持祖德留清白，冷落乡园愧弟昆；本分钳锤公透脱，更何了悟教诸孙。

扬州福国和尚至范赋二诗赠行

由此看来，郑板桥从家里的经验，对佛教当有所悉，印象深刻，亦颇尊重。

\vfill
\begin{BookNotes}
\BookNote{4.}{震华编纂，《兴化佛教通志》（台北，新文丰出版社公司，民国75年5月台一版)，页66。}
\end{BookNotes}

% 原书第 314 页（PDF 第 327 页）
\BookPage{314}{327}
\BookSubheading{（二）读书佛寺}

\begin{BookQuote}
《兴化佛教通志》提到郑板桥，「乞休归里，囊橐萧然，卖书画以自给，喜住
\end{BookQuote}

僧寺，寓焦山最久。」可见郑板桥在外，除仕官居衙署外，最佳的歇脚处就是僧寺，寺里有多余的禅房可读书、作画、会友、静思、居住，且有便餐可食，费用低廉，甚至无偿。因此，如谓佛寺是郑板桥的第二个家，实不为过。

佛寺跟郑板桥的结缘始于读书。当其年幼时，继母郝氏归宁时，携带同行，并向族叔郝梅岩请益受教，梅岩公当时借用庄内的净土庵东厢房开设学塾。从此揭启郑板桥与佛寺的不解之缘。因而郑板桥在兴化时也常跟书友同件于邻近寺庙念书。郑板桥曾追怀此一往事，曰：

\begin{BookQuote}
总角曾相聚，论文共一庭。诗歌春酒绿，风雨夜灯青。自折庭边柳，飘水上萍。寺钟天外远，何日得重听。〈寄怀兴化旧友》
\end{BookQuote}

康熙五十五年（1716），时年二十四岁的郑板桥考取秀才，为图家计，越明年，开馆设徒于真州江村（今江苏仪征）。然教书情形并不顺遂，颇有困顿失意之感，表示愿过渔隐生活。从此过着出游、卖画的生活·〈道情〉约在此后完稿，既想逃世，看破红尘；却又奋勉攻读应试经籍，内心充满徬徨与无奈。板桥自述一生共赴「秋闱」三次，第一次落榜时，第二次系恩科，卷子已被荐上，但因额满，仍未能中式，第三次总算能出头得意。8

由于郑板桥自惭为人父者，收入不佳，以致儿子营养不良，而惨遭爱子死亡之痛，逼使他鼓起勇气往科举之途发展，寄望借着「功名」能救「贫困」。在〈哭惇儿〉诗里，再三自我谴责，更见板桥对儿子之过逝归究于生活的穷苦所致。因此，为了第三次的秋试，郑板桥三十六岁之年，于雍正六年（1628）春，寓居扬州天

\vfill
\begin{BookNotes}
\BookNote{5.}{震华编纂，《兴化佛教通志》，页90。}
\BookNote{6.}{丁家桐，《绝世风流一郑燮传》（上海，上海人民出版社，2001年8月），页13。}
\BookNote{7.}{党明放，《郑板桥年谱》，页25-28。}
\BookNote{8.}{丁家桐，《绝世风流一郑燮传》，页74。}
\BookNote{9.}{黄泽雄，〈郑板桥生平及其思想之研究（上)〉，第一期《中台医专学报》（台中，中台医技专科学校，1983年)，页210。}
\end{BookNotes}

% 原书第 315 页（PDF 第 328 页）
\BookPage{315}{328}
\noindent 宁寺读书，研习制艺，潜心典籍，郑板桥之选择天宁寺，出于几项考虑：其一，天宁寺周边环境清静优雅，有利专注学习，避开家庭琐事的困扰。其二，寺里斋饭便宜，经济负担得起。其三，他爱与和尙交友阔谈，这里寺僧很多。在天宁寺期间，郑板桥常和陆白义、徐宗于诸砚友共同谈论诗文，情绪高昂，热血澎拜，有时兴之所至，还拔剑起舞，对酒当歌，发抒感怀。甚至还为了比赛记诵经文，彼此竞争默写。10

雍正十年（1732），郑板桥通过乡试，考上举人，再来更上层楼，就是三年一届的会试。雍正十三年(1735)，郑板桥为准备迎接次年的大考，就一个人搬到镇江焦山的佛寺去苦读。镇江在扬州对岸，依山傍水，山峰含黛，逶迤万状，焦山是江中的小岛，树木葱茏，其特色是寺庙甚多。板桥常住的庵庙是自然庵和别峰庵。一个人关在荒山古庵读书，如入囚笼，读书之余虽可听涛声，欣赏美景，但免不了浮想连翩，一方面求仕进，一方面想到政治社会的现实，一方面打算未来娶妾，一方面又想到身后的百年大事等，使得他得空就写信给墨弟，抒解苦闷无聊。同时也到各处寺庵闲逛，除跟和尚吟诗品茗，抵足长谈外，也在各寺庵留下对联墨迹。次年，即乾隆元年（1736)赴北京赶考，郑板桥被录取为二甲第八十八名进士。11简言之，郑板桥之能通过科考，寺庵提供的读书环境有所裨益，自事功不可没，也奠定郑板桥与佛教的密切关系。

\BookSubheading{（三）旅游居住}

古代文人出外远行旅游，离不开寺宿所提供的方便与照料。张潮《幽梦影》日：

\begin{BookQuote}
予尝谓二氏不可废，非袭夫大养济院之陈言也。盖名山胜景，我辈每思蹇裳就之，使非琳宫梵刹，则倦时无可驻足，饥时谁与授餐，忽有疾风暴雨，五大夫果真足恃乎？又或邱壑深邃，非一日可了，岂能露宿以待明日乎？虎豹蛇虺，能保其不人患乎？又或为士大夫所有，果能不问主人，任我之登陟凭吊而莫之
\end{BookQuote}

\vfill
\begin{BookNotes}
\BookNote{10.}{刘中建、林存阳，《郑板桥的狂怪人生》（北京，古籍出版社，2002年8月），页69-70。}
\BookNote{11.}{丁家桐，《绝世风流一郑燮传》，页89-94。}
\end{BookNotes}

% 原书第 316 页（PDF 第 329 页）
\BookPage{316}{329}
\begin{BookQuote}
禁乎？…即城市之内，与夫四达之衢，亦不可少此一种。客游可作居停，一也；长途可以稍憩，二也；夏之茗、冬之姜汤，复可以济役夫负戴之困，三也。凡此皆就事理言之，非二氏福报之说也。12
\end{BookQuote}

郑板桥的想法与做法，与此不谋而言，自谓：「酷嗜山水」，又说：「好山水，未能远迹，其所经历，亦不尽游趣。」〈板桥自叙〉既不能尽趣，当须再多游。故又曰；「板桥游历山水虽不多，亦不少；读书虽不多，亦不少；结交天下道人名士虽不多，亦不少。」〈板桥自序〉「虽不多」，事实上是很多的。试观他跑遍大江南北，旧地重游，就知道其好游的本性与程度。

郑板桥性爱山水，可能另有谋生的目的。试想，他经济欠佳，家人嗷嗷待哺，柴米不全，一家之主怎会有兴致去游览湖山？这需要川资，倘有余钱，不去养家糊口，怎么还去租借船马，在名山大川访僧寻道，终年不归。合理的解释，远出是为了利用种种关系沿门托钵，寻求机遇，说得好听一点，便是云游。13既然盘缠不足，又要远游，且交方外，基于经济的考量，最佳的居处就是寺庙，管游、管吃、管住、管友，真是一举数得，所以佛寺与郑板桥关系之深厚，概可想见。

康熙五十四年（1715）乙未秋九月，时年二十三岁的郑板桥以初生之犊，出游京师，寓北京瓮山漱云轩手书小楷欧阳修秋声赋〉轴并跋。跋曰：

\begin{BookQuote}
乙未九秋，山中寻菊，感黄叶之半零，望孤云而不返；残阳水面，渺渺寒涛，古寺山腰，凄凄晚碧，栖鸦欲定而犹惊，凉月虽升而未倾。偶翻欧赋，遂录是
\end{BookQuote}

篇。

瓮山及万寿山，郑板桥应住在山中古寺，漱云轩应为寺中一室，虽未知寺名，但文中已指证历历。14

京都岂有刚出茅庐的郑板桥所能容身，随即返乡，次年考取秀才，开馆授徒，生活艰苦，加上父亲弃世，惇儿夭折，更是辛酸凄惨。据其后来追忆，「学诗不

\vfill
\begin{BookNotes}
\BookNote{12.}{张心斋，《新校本幽梦影》（台北，天龙出版社，民国75年10月），页37-39。}
\BookNote{13.}{丁家桐，《绝世风流一郑燮传》，页46。}
\BookNote{14.}{党明放，《郑板桥年谱》，页24。}
\end{BookNotes}

% 原书第 317 页（PDF 第 330 页）
\BookPage{317}{330}
\noindent 成，去而学写，学写不成，去而学画，日卖百钱，以代耕稼，实救困贫，托名风雅」。〈署中示舍弟墨〉卖画的市场便是离兴化较近的扬州，其出道之时，约为三十一岁。这年（雍正元年）初春，郑板桥行价不佳，情绪低落，外游散心，至古代旧城海陵（今江苏泰陵），宿弥陀庵，开始与住持梅鉴上人订交。15

第二年，郑板桥出游江西，到过庐山。「初识上人在西江，庐山细瀑鸣秋窗。」。此上人指的是无方和尚，两人相谈甚欢。板桥有一幅竹是为无方画的，题曰：「最爱白方窗纸破，乱穿青影照禅床。」也许就是这一次会面画的，而这张床可能板桥就曾安寝过。接着亲临湖南洞庭湖，和洪觉范的潇湘八景诗，其中一首是〈烟寺晚钟〉，想当然又是投宿佛寺。16

郑板桥因只是秀才，缺乏名气，在扬州卖画，难以糊口，后得识尹会一、马曰琯，得其资助，于雍正三年（1725），再度游访京师。其目的一则为观赏北京的名胜古迹；次则试图寻求人生转折的机遇。这次的京师行，虽遇到佛门和尚、宫室子弟，仍四处碰壁，收效甚微，还是「衣裳检点不如归」的南下返乡。17惟客京期间，寓居慈仁寺，节省费用，且便与禅家尊宿深谈参机。18

雍正五年（1727），郑板桥有通州之旅，以其位在扬州东北，距离不远，抽空出游，也是住在不知名的佛寺，有诗为证：

\begin{BookQuote}
积雨空山草木多，山僧晨起所烟梦；崖前露出一块石，悄坐松阴似达摩。悬岩小阁碧梧桐，似有人声在半空；百叩铜环浑不应，松花满地午阴浓。
\end{BookQuote}

〈游白狼山〉

好一幅清爽寂寥的景致，既见「山僧晨起」，可见就在寺里过夜。19

游罢归来，依然如故，在扬州摆摊糊口，仍未能好转。几经挣扎，决定参加科考，以扩展生路，而于雍正六年（1728），一头栽进扬州天宁寺读书。四年后，

\vfill
\begin{BookNotes}
\BookNote{15.}{党明放，《郑板桥年谱》，页33。}
\BookNote{16.}{丁家桐，《绝世风流一郑燮传》，页46-47。}
\BookNote{17.}{刘中建、林存阳，《郑板桥的狂怪人生》，页56。}
\BookNote{18.}{党明放，《郑板桥年谱》，页42。}
\BookNote{19.}{党明放，《郑板桥年谱》，页47。}
\end{BookNotes}

% 原书第 318 页（PDF 第 331 页）
\BookPage{318}{331}
\noindent 赴金陵乡试，列榜举人。嗣后顺道游览南京名胜古迹，作〈念奴娇·金陵怀古〉词十二首，其中〈高座寺〉曰：

\begin{BookQuote}
暮云明灭，望破楼隐隐，院外青山千万叠，阶下流泉清浅，鸦噪松廊，鼠翻经匣，僧与孤云远，空梁蛇脱，旧巢无复归燕。可怜六代兴亡，生公宝志，绝不关恩怨。手种菩提心剑戟，先堕释迦轮转。青史讥弹，传灯笑柄，枉作骑墙汉。恒沙无量，人间劫数自短。〈高座寺》
\end{BookQuote}

在金陵，郑板桥游佛寺时，碰到一段失落的艳遇。诗曰：

\begin{BookQuote}
长干女儿年十四，春游偶过南朝寺。鬓发纤松拜佛迟，低头堕下金钗翠。
\end{BookQuote}

寺里游人最少年，闲行拾得翠花钿，送还不识谁家物，几嗅香风立怅然。〈长干女儿》

在南京期间，驻足多日，文献虽无记载住处，吾人推想，他在南京举目无亲，手头拮据，最好的落脚处还是佛寺，何况闲来无事，可跟和尙论对。接着从南京到杭州，先到钱塘江观潮，作〈观潮行〉与〈弄潮曲〉，后客居韬光庵，曾赋一诗，日：

\begin{BookQuote}
韬光古庵嵌山𪩘，北窗直吸余杭县；葛洪小儿峰岭低，南屏一片排秋扇。…山中老僧貌奇古，十年不踏西冷土；厌听湖中歌吹声，肯来伺候衙门鼓？…我已无家不愿归，请来了此前生果。〈韬光〉
\end{BookQuote}

在韬光庵中为方丈松岳上人作画，并题五绝一首，同时也写家书〈雍正十年杭州韬光庵中寄舍弟墨〉，说些自己的想法。在杭州，接到考中举人的喜讯后，写下〈得南阐捷音〉诗，描述内心复杂的感受。20

郑板桥中举，应于次年参加会试，然这一年，叔父省庵公亡逝，因而没有应试，大概与「居忧」有关，因板桥父母已故去，只有这个叔父，且省庵公待板桥甚好，所以板桥为之执「居忧」礼，也是情理中事。加上「身患大疮」，严重到「浑身动弹不得」，更不良于行。既然不参加会试，时间就很充裕，于是第二次出游海陵，仍跟前次相同，寄住在弥陀庵的寺庙中，方丈梅鉴和尚与板桥早在雍正元年

\vfill
\begin{BookNotes}
\BookNote{20.}{党明放，《郑板桥年谱》，页62-63。}
\end{BookNotes}

% 原书第 319 页（PDF 第 332 页）
\BookPage{319}{332}
\noindent （1723）就有交往，此次重逢，很是相得投契。上人不修边幅，酷爱诗文，时常请教板桥题诗写字，惟其如此，郑板桥为他做两首诗，即〈别梅鉴上人〉、〈赠梅鉴上人〉。重访海陵，均住弥陀庵，可见二人情谊之身，不言可喻。21

海陵归后，仍在扬州交友、卖画，不再远游。雍正十三年更继续未竟之业，全力以赴，准备会考，特赴镇江焦山别峰庵闭门苦读，其间虽曾到处走走，调剂身心，然亦流转于各佛寺间，朝夕相见的大都是和尚。乾隆元年（1736）二月到北京参加礼部会试，中选二甲进士，满心欢喜，遂赋〈秋葵石笋图〉以自贺。但郑板桥并未立即返乡，而是逗留京师，其目的期待求得任官派职，或于庶常馆磨练，曾曰：「余本书生，初志望得一京官，聊为祖父争气。」然而郑板桥丑陋的容貌，横溢惊座的才气，加上孤傲颠狂的性格，都是进入仕途的大忌。一年后，终无所获，怏怏地南归回乡。22

在京期间，待官无望，虽甚悻然，但长住佛寺，与和尚交往，既可满足旅游愿望，且能畅谈适性，应会稍释情怀。至瓮山跟无方上人叙旧，作〈赠瓮山无方上人二首〉诗，后又作〈瓮山示无方上人〉诗，曰：

\begin{BookQuote}
松梢鴈影度清秋，云淡山空古寺幽。蟋蟀乱鸣黄叶径，瓜棚半倒夕阳楼。客来招饮欣同出，僧去烹茶又小留。寄语长安车马道，观鱼濠上是天游。
\end{BookQuote}

〈瓮山示无方上人》

郑板桥也曾到香山卧佛寺访青崖和尚，作〈访青崖和尙·和壁间晴岚学士虚亭侍读原韵〉诗，后又赋〈寄青崖和尚〉诗，曰：

\begin{BookQuote}
山中卧佛何时起，寺里樱桃此日红。骤雨忽添崖下水，泉声都作晚来风。紫衣郑重君恩在，御墨淋漓象教崇。透脱儒书千万轴，遂令禅事得真空。
\end{BookQuote}

〈寄青崖和尚〉

接着出游法海寺见仁公，作〈法海寺访仁公〉诗，又作〈同起林上人重访仁公〉诗，曰：

\vfill
\begin{BookNotes}
\BookNote{21.}{刘中建、林存阳，《郑板桥的狂怪人生》，页58-59。}
\BookNote{22.}{刘中建、林存阳，《郑板桥的狂怪人生》，页76-78。}
\end{BookNotes}

% 原书第 320 页（PDF 第 333 页）
\BookPage{320}{333}
\begin{BookQuote}
几日不相见，作诗盈一囊，立残云外露，销尽定中香。雨歇四天碧，风高秋稼黄。可应歌「击壤」，更为继陶唐。〈同起林上人重访仁公》
\end{BookQuote}

后来又作〈山中夜坐再陪起林上人〉诗，提到「晨起望诸山，烟岚漭涨塞」。既然「夜坐」、「晨起」，可见郑板桥与上人同住禅房。23除佛寺外，道观也曾进住，兹以〈宿光明殿赠娄真人〉诗为证。试想郑板桥在北京的日子里，卖画不佳，盘缠短少，又人地生疏，缺乏外援，寺庙自是最好的居处。只是文集中无法详细叙述，而佛寺有新科进士郑板桥到访，也是增光添彩的事。

郑板桥由京师回到扬州后，直到乾隆六年（1741），他都在兴化与扬州度过的。中进士前，板桥在扬州的住所经常变更，但大都在寺庙里的。乾隆六年九月，郑板桥奉吏部之召，再次来到京师，重逢雍正三年在北京结识且给他许多帮助的皇叔慎郡主允禧，也许经其居间说项，板桥不久就被掌理山东范县。到北京后，依然喜欢在佛寺挂单。如赴京都前，作〈逢客入都寄勖宗上人口号〉诗，先行传话，虽是问候，实则预为借宿。曰：

\begin{BookQuote}
汝到京师必到山，山之西麓有禅关，为言九月吾来住，检点白云房半闲。〈逢客入都寄勖宗上人口号》
\end{BookQuote}

前此熟稔的和尚，也借此拜访小住，以享受出游的闲趣。如致意青崖老人，诗日：

\begin{BookQuote}
一夜西风雪满山，老僧留客不开关；银沙万里无来迹。犬吠一声村落间。〈山中卧雪呈青崖老人》
\end{BookQuote}

因此，郑板桥一生中的出外旅游或参加科考，所到之处，尽量住在佛寺，无论城市古刹，荒山野寺，他都宾至如归，安之若饴。十年知县期间，或外出或南归，亦复如是。及至辞官返乡卖画，名气声望与艺术成就更是提高，虽可住更好馆舍，但郑板桥仍习惯寺院，取其宽敞清雅。如回到扬州，郑板桥住在城北竹林寺，城中爱好书画斯文的年轻人，都争相上门请益，来自四面八方的文人学士也经常去那里

\vfill
\begin{BookNotes}
\BookNote{23.}{党明放，《郑板桥年谱》，页84-86。}
\end{BookNotes}

% 原书第 321 页（PDF 第 334 页）
\BookPage{321}{334}
\noindent 研讨艺文方面的问题。24

\BookSubheading{（四）方外订交}

文人向来都跟佛教有些渊源，当然包括与和尚的互动往来，其中最被人称道的是郑板桥。慧禅独具见识的编撰《名人与佛教》上、下两册，选取历代名人居士共一百八十四位，叙述他们对佛教的观点与素养，及跟佛教间的轶事趣闻；更重要的是，每个人都用一、二句话作刊头标语，以收画龙点睛之效。如郑板桥为「平生不糊涂，好与僧侣游」，就非常贴切。25因前句乃「难得糊涂」转化而来，至于「好与僧侣游」则是事实。一百八十四位名人中，特别强调郑板桥的这个特点。换言之，在作者的眼光里，历代文人与方外的交往，都不如郑板桥来的广阔与热络。

郑板桥之所以有如此多的方外之交，亦因他性爱山水，即使在穷困落魄时，依然游兴不减。出于自身的喜好，同时考虑到经济原因，板桥常到寺庙投宿，在此情况下，出家人都乐于提供方便，板桥因而交到不少僧道朋友。26何况郑板桥天真率直，深通佛里，书画绝佳，诙谐笑骂，颇得僧侣们的喜爱。

几乎每本郑板桥的传记都会记述其僧侣友人，当中最详细的是金宝秋大作〈郑板桥方外友考略〉，文曰：

\begin{BookQuote}
据有关资料，板桥之方外友竟达三十余位之多，其诗章唱和、信函往来、书画酬赠，满载着他们之间有交游的深笃情谊，记录着他们之间心灵的相知相契。这三十几位和尚，大都非一般念经念佛之徒，如无方、青崖、仁公、莲峰、弘量、巨潭、福国、焚虚等，有为朝廷所重者、有为普门龙象者、有雅擅翰墨者、有德行高洁者，于一时一方均声名卓着，板桥喜与之游，岂能在思想上、艺术上没有其影响乎？27
\end{BookQuote}

\vfill
\begin{BookNotes}
\BookNote{24.}{引自杨樱林、黄幼钧编着，《中国书画名家画语图解—郑板桥》（北京，中国人民大学出版社，2010年2月），页103。}
\BookNote{25.}{慧禅主编，《名人与佛教》，上册（上海，上海人民出版社，2008年10月），页250-252。}
\BookNote{26.}{刘中建、林存阳，《郑板桥的狂怪人生》，页180。}
\BookNote{27.}{金实秋，〈郑板桥方外友考略〉，《郑板桥与佛教禅宗》，页1。}
\end{BookNotes}

% 原书第 322 页（PDF 第 335 页）
\BookPage{322}{335}
该文特别简述郑板桥与其中二十二位释子间的往来情形。进而作「郑板桥方外友简表」，列出三十五位僧侣的法号、当时身分、交往时间与地点、资料出处等项，清楚明确，极便理解。28此外，作者有感于无方上人与郑板桥接触较深，交游最密，历时最长的一位，现存作品中（疑义的板桥作品除外），提及无方凡七处，是板桥作品中极少见的，因而另行撰述〈郑板桥与无方〉一文，以记其往还关系。29

党明放的佳构《郑板桥年谱》内容丰富详实，为研究郑板桥的绝佳参考，其附三郑板桥僧友道友人名录〉，提到郑板桥的僧友共三十七位。僧友包括无方上人、恒彻上人、福国上人、松岳上人、梅鉴上人、宏量上人、碧崖上人、博也上人、巨潭上人、松风尙人、勖宗上人、慧圆上人、慧鉴上人、侣松上人、成岳上人、宏熔上人、莲若上人、起林上人、佛上人、青崖和尚、啸江和尚、仁公和尚、拙公和尚、彼公和尚、永公和尚、药根和尚、偈船和尚、慧通和尚、慧如和尚、克恭和尚、六安和尚、莲峰上人、笠矶上人、立思上人、静山上人、复显上人、瑟虚上人。这些和尚与郑板桥的关系散见书内各页，为省篇幅，不作说明。30

上述诸上人外，郑板桥在文集中还常提到一些不知名的老僧、山僧，可见他对和尚们深具气味相投，广结善缘的情结。

\BookSubheading{三、郑板桥的无佛之佛}

郑板桥博通儒释道，兼构三家思想，尽管与佛教的关系深远流长，笃厚诚挚，但终究非我佛中人，不列释徒之林。

\BookSubheading{（一）归本儒家}

中国文化蕴涵儒释道三家，历代文人都多所研习，并受其影响、惟各家程度不同而已。宋明理学兴起后，儒家居主流地位，郑板桥在这个大环境下，虽有佛道

\vfill
\begin{BookNotes}
\BookNote{28.}{金实秋，〈郑板桥与佛教禅宗〉，《郑板桥与佛教禅宗》，页105-107。}
\BookNote{29.}{金实秋，〈郑板桥与无方〉，《郑板桥与佛教禅宗》，页68-71。}
\BookNote{30.}{党明放，《郑板桥年谱》，页460-462。}
\end{BookNotes}

% 原书第 323 页（PDF 第 336 页）
\BookPage{323}{336}
\noindent 的色彩，但基本上仍归本儒家。故论者谓郑板桥「日放高谈藏否人物，以是得狂名」。儒哲亦言「狂者进取」，狂固然是郑板桥的特点，但其心魂深处仍体现儒家的本质，敬仰圣贤的思想，表现仁者爱人的情怀。31

据文献所示，郑板桥可说是儒学世家，数代传承。曾祖名新万，字长卿，为邑庠生。祖父名湜，曾任儒官。父名之本，字立庵，号梦阳，品学兼优，家居授徒，皆有成就。叔之标，字省庵，生子墨，字五侨，庠生。外祖父汪翊文，奇才博学，隐居不仕，继母族叔郝梅岩亦为名儒。可见板桥当时的家世，虽渐趋没落，但称的上学脉相沿，儒门世家，文章品行，历历赓续。32

为了解郑板桥的儒家思想，先从其读书说起。郑板桥在扬州天宁寺时，常与学伴诵读之余，默写四书五经。曰：

\begin{BookQuote}
四书五经自家又未尝时刻而稍忘，无他，当忘者不容不忘，不当忘者不容不不忘耳。戊申之春，读书天宁寺，咕哔之暇，戏同陆、徐诸砚友赛经生熟。市坊间印格，日默三五纸，或一二纸，或七、八、十余纸；或兴之所至，间可三二十纸，不两月而竣工。虽字有真草讹减之不齐，而语句之间，实无毫厘错谬。固诵读之勤，亦刻苦之验也。〈四字书真迹序〉
\end{BookQuote}

郑板桥给其弟墨的书，强调如何读书。曰：

\begin{BookQuote}
四书之上有六经，六经之下有左、史、庄、骚、贾、董策略，诸葛表章，韩文杜诗而已，只此数书，终身读不尽，终身受用不尽。至如二十一史，书一代之事，必不可废。…六经之文，至矣尽矣，而又有至之至者：浑沦磅礴，阔大精微，却是家常日用。〈焦山别峰庵雨中无事书寄舍弟墨》
\end{BookQuote}

这是郑板桥的读书心得，儒家典籍浩如烟海，板桥沉浸良深，由博返约，取精用宏，提出要在「终身受用不尽」的书上做工夫，足见其早年读书应举之艰苦。信中所开列的书目，可以窥知其选录的标准，首先要求的是以民为本，有益于社稷苍

\vfill
\begin{BookNotes}
\BookNote{31.}{郑德开，〈郑板桥与儒释道〉，《楚雄师范学院学报》，第22卷11期（楚雄师范学院，2007年11月），页39。}
\BookNote{32.}{王建生，《郑板桥研究》（台北，文津出版社，1999年8月），页7-8。}
\end{BookNotes}

% 原书第 324 页（PDF 第 337 页）
\BookPage{324}{337}
\noindent 生，这是标准的儒家仁政思想。33

次从思想而论，儒家主张内圣外王，内圣则忠恕、仁爱、修德。郑板桥给墨弟书曰：

\begin{BookQuote}
余五十二岁始得一子，岂有不爱之理！然爱之必以其道，虽嬉戏顽耍，务令忠厚悱恻，毋为刻急也。平生最不喜笼中养鸟，我图娱悦，彼在囚牢，何情何理，而必屈物之性以适吾性乎？我不在家，儿子便是你管束。要须长其忠厚之情，驱其残忍之性，不得以为犹子而姑纵惜也。〈潍县署中与舍弟墨第二书》
\end{BookQuote}

\%

另一家书云：

\begin{BookQuote}
或曰：此等宅居甚适，只是怕盗贼。不知盗贼亦穷民耳，开门延入，商量分惠，有甚么便拿甚么去；若一无所有，便王献之青毡，亦可携取质百钱救急也。〈范县署中寄舍弟墨第二书
\end{BookQuote}

像这种言论，尚有多处，可见郑板桥深寓忠厚心怀，仁民爱物、推己及人的精神，完全符合儒家要义。

至于外王就是做善事、行德政。郑板桥为秀才时，「捡家中旧书簏，得前代家奴契券，即于灯下焚去，置不返诸其人」（〈雍正十年杭州韬光庵中寄舍弟墨〉）。如此存心，可谓善事，类此义举，终生俸行。故其家书提及「切望我弟多行善，彼苍虽属莽莽，惟积德之报，屡试不爽」。再者，郑板桥为官清廉公正，打击劣绅恶商，关心民间疾苦，抚恤灾难百姓，甚至不惜抵触法制，自行开仓赈贷，救济灾民。此一爱民表现，就是具体的儒家实践。

复次，论及尊孔思想，以郑板桥之狂怪任性，却非常宗奉孔子，认为孔子是天生圣人，后无来者。曰：

\begin{BookQuote}
天生圣人亦屡矣，未尝生孔子也。及生孔子，天地亦气为之竭而力为之衰，更不复能生圣人。〈潍县署中与舍弟墨第二书·书后又一纸》
\end{BookQuote}

尧与舜都是中国古代圣贤名君，同样受到景仰，但郑板桥就崇尧降舜，即受孔

\vfill
\begin{BookNotes}
\BookNote{33.}{郑德开，〈郑板桥与儒释道〉，页38。}
\end{BookNotes}

% 原书第 325 页（PDF 第 338 页）
\BookPage{325}{338}
\noindent 子影响所致。曰：尝论尧舜不是一样，尧为最，舜次之。人咸惊讶。其实有至理焉。孔子曰：「大哉，尧之为君，惟天为大，惟尧则之。」孔子从未尝以天许人，亦未尝以大许人，惟称尧不遗余力，意中口中，却是有一无二之象。…尧乎，尧乎！此其所以为天也乎！非天之薄于大舜而不予以后也，其道已尽，其数已穷，更无从蕴而再发耳。〈潍县署中与舍弟墨第二书·书后又一纸。〉

郑板桥亟为赞赏孔子的读书，虽然做到圣人，也要困勉力学，曰：读书以过目成诵为能，最是不济事。眼中了了，心下匆匆，方寸无多，往来应接不暇，如看场中美色，一眼即过，与我何与也。千古过目成诵，孰有如孔子者乎？读《易》至韦编三绝，不知翻阅过几千百遍来，微言精义，愈探愈出，愈研愈入，愈往而不知其所穷。虽生知安行之圣，不废困勉下学之功也。〈潍县署中寄舍弟墨第一书》

孔子不仅会读书，然因汗牛充栋，更懂得删书，并以秦始皇烧书之误相对比，以判高下，而有仙凡之别，曰：秦始皇烧书，孔子亦烧书。删书断自唐、虞，则唐、虞以前，孔子得而烧之矣。诗三千篇，存三百十一篇，则二千六百八十九篇，孔子亦得而烧之矣。孔子烧其可烧，故灰灭而无所复存，而存者为经，身尊道隆，为天下后世法。始皇虎狼其心，蜂虿其性，烧经灭圣，欲剜天眼而浊人心，故身死宗亡国灭，而遗经复出。始皇之烧，正不如孔子之烧也。〈焦山别峰庵雨中无事书寄舍弟墨》

更有趣的是，郑板桥藉金陵圣庙倾倒，以孔子之圣批判冬烘秀才，曰：噫！圣庙工程何等坚固，自建筑至今，已阅数千年，经历之暴风急雨，不胜以偻指计，而至今宫墙间苔痕尺许厚，不类有倾倒重砌象，何独不耐今届之风雨而顷颓乎？无他，盖因金陵城中龌龊秀才满坑满谷，现任教喻，亦属胸中绝无点墨者。斯文扫地，辱没圣门，孔子岂容若辈列门墙，故特毁墙以示驱逐之意，殆其然乎？则予亦复何言哉！复何言哉！〈寄墨弟自焦山发》

从上看来，郑板桥的读书大都儒家经典为多，虽亦旁即史籍、文集，却鲜有提


% 原书第 326 页（PDF 第 339 页）
\BookPage{326}{339}
\noindent 及佛典。而思想行为，称得上是狂怪，然终究未脱离儒家理念。尤其是崇敬孔子的程度，与理学家亦无轩轻。至于郑板桥的诗文中，虽多处间涉佛教，且与方外的交情亦过从甚密，但并未皈依佛门，且少有讨论佛教的本质问题。因此，尽管郑板桥兼修儒释道，但核心价值仍是儒而非佛。简言之，郑板桥在意念上，非佛门中人，心中并无真佛，却又佛缘殊胜，故可谓无佛之佛。

\BookSubheading{（二）民间信仰}

郑板桥之未成为优婆塞或在家念佛，除认同儒家外，其对民间信仰的态度，也可得到理解。基本上，郑板桥在传统儒家思想的薰陶下，接纳道家、阴阳家与佛家的元素，就宗教信仰而言，大致仍属传统社会的民间信仰，只是多些儒家的忠恕仁爱，少点迷信乱神。

「天道」观念，中国古已有之，其意涵虽称精奥深博，却也含糊笼统，但民众确信不疑，郑板桥也不例外。曰：

\begin{BookQuote}
王侯将相，岂有种乎！而一二失路名家，落魄贵胄，借祖宗以欺人，述先代而自大。辄曰：彼何人也，反在霄汉；我何人也，反在泥涂。天道不可凭，人事不可问。嗟乎！不知此正所谓天道人事也。天道福善祸淫，彼善而富贵，尔淫而贫贱，理也，庸何伤？天道循环倚伏，彼祖宗贫贱，今当富贵，尔祖宗富贵，今当贫贱，理也，又何伤？天道如此，人事即在其中矣。〈雍正十年杭州韬光
\end{BookQuote}

庵中寄舍弟墨〉

可见郑板桥以为冥冥之中，有一种「天道」，循环倚伏·操纵着人间的富贵贫贱，所以，人应该乐天知命。但是，他却也不是全然悲观的宿命论者，他同时又表示所谓的「天道」，有时又是凭借着人事而来。天道福善祸淫，所以，人不可以不问人事，但是尽人事还得听天命，这是承袭孔门乐天知命的说法。能如此，则能不怨天、不求人，也不再长叹「天道不可凭，人事不可问」，人人奋力自强，则「将目无种」，谁不可以为王侯？谁不可以为尧舜？34

\vfill
\begin{BookNotes}
\BookNote{34.}{廖玉惠选注，《一竿烟雨一郑板桥诗词文选》（台北，时报文化出版公司，2000年7月），页252。}
\end{BookNotes}

% 原书第 327 页（PDF 第 340 页）
\BookPage{327}{340}
郑板桥除相信天道，还受到阴阳五行的影响。曰：

\begin{BookQuote}
夫天地生物，化育劬劳，一蚁一虫，皆本阴阳五行之气𬘡而出。上帝亦心心爱念，而万物之性人为贵，吾辈竟不能体天之心以为心，万物将何所托命乎？蛇蚖蜈蚣、豺狼虎豹，虫之最毒者也，然天既生之，我何得而杀之？若必欲尽杀，天地又何必生。〈潍县署中与舍弟墨第二书〉
\end{BookQuote}

郑板桥的这种观念，与道家亦有相通之处。他说：「将但时，睡梦初醒，尚辗转在被，听一片啁啾，如云门威池之奏，及披衣而起，颓面漱口啜茗，见其扬翚振彩，倏往倏来，目不暇给，固非一笼一羽之乐而已。大率平生乐处，欲以天地为囿，江汉为池，各适其天，斯为大快。」所谓「天地为囿，江汉为地，各适其天」之思想，正与庄子所云「天地与我并生，而万物与我为一」的观念相似，又庄子云:「故曰天在内，人在外。德在乎天，知天人之行，本乎天，位乎得。…牛马四足，是谓天！落马首，穿牛鼻，是谓人。故曰无以人灭天，无以故灭。」之「天人合一」的理论，正是郑板桥以「天地为囿，江汉为池，各适其天」的思想。35

既然接受所谓天道与人道的说法，郑板桥免不了要碰上风水算命等问题，对这些堪舆术数，他则处在信与不信，时信时不信之间，很难跟民间信仰割舍。就不信的情况来说，中国是农业社会，靠天吃饭，从事耕种，需水灌漑，如遇干旱，就歉收灾荒，此时官民要向天神祈雨，解决危机。此项行为，郑板桥认为无效，加以讥嘲。曰：

\begin{BookQuote}
祈雨能下雨，高山可种田；吃斋能行善，骡马上西天。求神能通灵，谁人还遭难；有钱能买命，天子万万年。〈嘲祈雨〉
\end{BookQuote}

传统民间向来重视风水，各家说法，模棱两可，各有差别，甚至冰炭不容，郑板桥则认为只要心地善良，葬那里都好，故曰：

\begin{BookQuote}
总之，风水之术，原属渺茫难凭，堪舆家画蛇添足，非富即贵，打动人心，借以赚钱，「寅葬卯发」，我则但闻其说，未见其事也。愚意选择墓地，不必望大富大贵，要子孙如何发达，只要地位清佳，地土温
\end{BookQuote}

\vfill
\begin{BookNotes}
\BookNote{35.}{黄泽雄，〈郑板桥生平及其思想之研究（下）〉，《中台医专学报》，第二期（台中，中台医技专校，民国73年)，页176-177。}
\end{BookNotes}

% 原书第 328 页（PDF 第 341 页）
\BookPage{328}{341}
\noindent 腴，墓外有气势可观，穴中无蚊无水，埋葬其间，先灵永妥，不妨自家，不碍别人，便是好风水耳。如必欲一门发达，代代兴旺，雄武如秦之始皇帝，大可拣选天下最好之地而葬之，使子孙万世，绵延不绝，何至沙丘身殒，一传而即斩灭宗社乎？君子于此，亦可警悟。吾扬有谚曰：「心地好，阴地超」。心地纯厚，阴地不好而自好，吾何必信惑术者，改葬易向，使先灵不安于窀穸中乎，足下明达，当能体会斯言。〈答汪希林〉

这段话是对朋友说的，郑板桥更剑及履及，自身实践，买墓地不重风水。曰：郝家庄有墓田一块，价十二两，先君曾欲买置，因有无主孤坟一座，必须刨去。先君曰：「嗟乎！岂有掘人之冢以自立其冢者乎！」遂去之。但吾家不买，必有他人买者，此冢仍然不保。吾意欲致书郝表弟，问此地下落，若未售，则封去十二金，买以葬吾夫妇。即留此孤坟，以为牛眠一伴，刻石示子孙，永永不废，岂非先君忠厚之义而又深之乎！夫堪舆家言，亦何足信。吾辈存心，须刻刻去浇存厚，虽有恶风水，必变为善地，此理断可信也。〈焦山双峰阁寄舍弟墨〉

如此看来，郑板桥似乎反对风水，不信鬼神，但事实并不尽然。再就相信的情况来说，他还是承认风水的好处，曰：刹院寺祖坟，是东门一枝大家公共的，我因葬父母无地，遂葬其傍。得风水力，成进士，作宦数年无恙。是众人之富贵福泽，我一人夺之也，于心安乎不安乎!〈范县署中寄舍弟墨〉

有趣的是，郑板桥与风水先生耿仲南友善，建宅、选墓皆请他勘察决定，并深信不疑。曰：曾托堪舆耿仲南为余家相宅，谓门向偏东，远挹山光，近临水秀，不出显宦，必出名士。惟宅后交叉两路，旺气被行人踏破，丁口难期旺盛耳。其时余仅一酸秀才，未及十年，果成进士。耿氏之言，洵有微也。

郑板桥虽曾说:「夫堪舆家，亦何足信。」但此为说说而已，骨子里还是笃


% 原书第 329 页（PDF 第 342 页）
\BookPage{329}{342}
\noindent 信不疑的，所谓的「何足信」者，乃已知其不灵之故耳。36所以郑板桥并非不信风水，只是要找到他所信任且满意的堪舆家。

算命也碰到同样状况。有位盲词人陈孟周曾替郑板桥算过八字，可能将命测得太好了，致使郑板桥罢官后，「乃知孟周之言亦不灵也」。但又相信另一位算命师，他给郑墨的信即曰：

\begin{BookQuote}
杨四衙儿子命理甚精，比俗流欲高数等，谓我这五年是晦气，…我六十五岁方行大运，与前不同，常为内京官、掌生杀。湖州太守命学精，谓我六十五后生子，扬名发财，其命章带与你看。若果如此，吾弟可无忧窘隘矣。个个算命皆如此说，而杨、李二公谈得最为亲切有理，咬牙顿口不差。可与太太、两嫂子并大女、二女说也。37
\end{BookQuote}

由此看来，郑板桥对风水、算命颇有兴趣，亦有些心得，找过数家相比较，说得准的就相信探纳，否则不予探任。郑板桥的这种表现，不正是中国民间信仰具体而微的缩影吗？

如前所述，郑板桥在〈嘲祈雨〉诗，表示向神祈雨是没有用的，仿佛神是无法通灵显赫，但在自编文集后，不愿被人添加窜改，又有不同说辞，曰：

\begin{BookQuote}
板桥诗刻止于此矣，死后如有托名翻板，将平日无聊应酬之作，改窜烂入，吾必为厉鬼以击其脑。〈后刻诗序》
\end{BookQuote}

郑板桥在不同的两首诗，对神鬼做分歧的解读，似乎有违常态，但如放在民间信仰的思维中，因时空变化，处境有别时，表现出互异的心境与情绪，应是料想得到的事。在生活事物上，应无意外可言。

当人年龄愈长，阅历愈多，感慨相对增加，宗教情怀也逐渐加深，尤其是遇到重大事故，情况更为明显。郑板桥五十七岁时，唯一的儿子与挚友高凤翰相继逝世，令他深感痛心，无限惆怅。次年，撰〈文昌祠记〉，起首曰:「文云乎哉！行云乎哉！神云乎哉！修其文，懿其行，祀其神，斯得之矣。」可见已走进神的境

\vfill
\begin{BookNotes}
\BookNote{36.}{金实秋，〈郑板桥与道士〉，《郑板桥与佛教禅宗》，页61。}
\BookNote{37.}{引自金实秋，〈郑板桥与道士〉，《郑板桥与佛教禅宗》，页60-61。}
\end{BookNotes}

% 原书第 330 页（PDF 第 343 页）
\BookPage{330}{343}
\noindent 界。六十岁时，把全部精神寄托在宗教的性灵上，捐银五十两，倡修潍县城北玉清宫，并另述〈城隍庙碑记〉，表达其对宗教的看法。38

虽然郑板桥跟佛教寺僧的关系远较宫观道人来的密切，替佛教所写的诗书联对也超过道教甚多，但大致看来，郑板桥的宗教情怀与精神心灵是依附在以儒家为主的的民间信仰上。

\BookSubheading{（三）无佛之佛}

从郑板桥率真、狂妄、怪诞、殊异的性格，很难成为皈依三宝，一心向佛的信徒，却非常适合跟出家人闲话谈禅，论说佛理，因而与方外订交，非常投缘，情恳意深。由于在儒家思想的社会，加上世代儒学传家，使得郑板桥自然地接受儒家薰陶，崇仰孔子。而且传统民间信仰，风水算命的习俗，多少也影响了郑板桥的想法，尽管郑板桥和佛教、佛学、佛寺、佛僧有长期深厚的关系，但没有动摇他的基本思想与核心观念，却能丰富思维意念，充实生活内涵，增广学问素养。因此，郑板桥是在儒家的基础，吸纳佛家，以建立其特有的人生哲学与生命情怀。

金宝秋也抱持同样的见解，认为郑板桥深受佛教禅家的影响是事实，值得吾人思考与研究。但是，应当指出的是，郑板桥对佛门禅宗，只止于「禅悦」而已，并非是虔诚的信仰崇拜者；他之所以喜欢与禅宗尊宿交游，除去思想相通、情趣相类之原因外，也是向社会表示其孤傲清高、绝尘脱俗。再者，板桥的思想是复杂的，多变的，儒、道、释三种中国传统文化在他的作品上都有所反映、适时体现。我们没有必要、也很难如贴标签那样，将三家的影响一一区分得清清楚楚；况且，三家本身也在不断地相互渗透、吸收而交融。综观板桥一生，儒家正统思想始终居于主导地位，他还不能突破他生活的那个时代，他对「乾隆东封书画史」尚且念念不忘，对「康熙秀才、雍正举人、乾隆进士」那样津津乐道，对「我亦终葵成进士，相随丹桂状元郎」如此洋洋自得，就是明证。39

雍正三年（1725)，郑板桥北游燕京，因「日放言高论，臧否人物，无所忌

\vfill
\begin{BookNotes}
\BookNote{38.}{黄泽雄，〈郑板桥生平及其思想之研究（上）〉，《中台医专学报》，第一期，页217。}
\BookNote{39.}{金实秋，〈郑板桥与佛教禅宗〉，《郑板桥与佛教禅宗》，页104。}
\end{BookNotes}

% 原书第 331 页（PDF 第 344 页）
\BookPage{331}{344}
\noindent 讳」。加上性格上的「自负」、「好大言」、「谩骂无择」，就写〈燕京杂诗〉。

\begin{BookQuote}
不烧铅汞不逃禅，不爱乌纱不要钱。但愿清秋长夏日，江湖常放米家船。
\end{BookQuote}

〈燕京杂诗〉之一

这首诗明显看出郑板桥的四不，即不道、不佛、不作官、不要钱，其非佛家的心态，早就表露无遗。

郑板桥在另一诗中，表示深沉的感慨。曰：

\begin{BookQuote}
古来开国多圣明，历尽艰难身百战；一朝勘定称至尊，承明殿上头毛变。安期枣尽还瘦羸，赤松皇帝坟垒垒；学先学佛空尔为，昼苦短，西日飞。
\end{BookQuote}

〈昼苦短》

此诗说明开国帝王，历尽艰辛，征战沙场，终于功成名就，后来白发年老，「赤松皇帝坟垒垒」意谓古来神仙，仍不免一死，难怪郑板桥叹说「学仙学佛空尔为」，仙佛亦为空话，无济于事。40

\begin{BookQuote}
〈偶然作〉也是郑板桥心态的照映，曰：英雄何必读书史，直摅血性为文章；不仙不佛不圣贤，笔墨之外有主张。纵横议论析时事，如医疗疾进药方…。〈偶然作》
\end{BookQuote}

这首诗提示英雄文章不必顾虑仙佛，贤圣之说的基本原则，所谓「笔墨之外有主张」，就是要有真正理念，即明确的修齐治平之实用论。因此，他说不管你如何纵横议论，分析时事，都必须像医家进药方一样，针对时弊，其最终目的，还是要能治病，可为社稷民生解决问题，绝对不能如仙佛与有些贤圣，徒托空言。诗的后段以三种层次来区别文章的优劣，而大体以儒家「征圣、宗经、明道」等实用文章观为鹄的。所以他作品的题材多半出于现实，来自生活，言之有物，感触而发，使人领悟到郑板桥强力、丰富、诚挚的「民胞物与」情怀。41此一笔墨主张，显然不是佛家的立论所在。

\vfill
\begin{BookNotes}
\BookNote{40.}{王连生，《郑板桥研究》，页198-199。}
\BookNote{41.}{廖玉惠选注，《一竿烟雨一郑板桥诗词文选》，页25-26。}
\end{BookNotes}

% 原书第 332 页（PDF 第 345 页）
\BookPage{332}{345}
\noindent 雍正元年（1723）郑板桥三十一岁时，独子逝世，悲伤欲绝，作〈哭惇儿诗，曰：蜡烛烧残尚有灰，纸钱飘去作尘埃。浮图似有三世生说，未了前因好再来。〈哭惇儿〉之五

诗虽曾提到浮图三世，因在丧子之痛的情况下，用词难免比较会有宗教意识，何况另加「似有」，表示并未尽信。因此，郑板桥对佛教信仰的问题，心中自有定见，非常明确。曰：

僧人偏满天下，不是西域送来的。即吾中国之父兄子弟，穷而无归，入而难返者也。削去头发便是他，留起头发还是我。〈焦山读书寄舍弟墨》削发的是出家人，留发的则非佛教徒，这个「我」就是郑板桥自喻，所以郑板桥不信佛教，是很清楚的事。向来文人论郑板桥以诗书画为主，蒋士铨虽不例外，有所详论，然能特具灼见，看出他不是一般常人，但也走不进佛门殿堂。诗曰：

未识顽仙郑板桥，非人非佛亦非妖；晚摹「痤鹤」兼山谷，别开临池路一条。42

郑板桥固然没有皈依佛教，却十分向往佛教的悠闲情境，诗词中多处提及，兹录一首，以示其意义。曰：

茅庵欹欲倒，倩老树撑扶，白云环绕。林深无客到，有涧底鸣泉，谷中幽鸟。清风来扫，扫落叶尽归炉灶。好闭门煨芋挑灯，灯尽芋香天晓。非矫，也亲贵胄，也踏红尘，终归霞表。残衫破被，补不彻，缝不了。比世人少却几茎头发，省得许多烦恼。向佛前烧柱香儿，闲眠一觉。〈瑞鹤仙·僧家》

身处这样一个老树白云、澄澈幽雅的环境，寺庭明净，挑灯看书，真是人间仙域。但也不必自命清高，照样可以踏入红尘，只是终归霞表。残衫破袖，既不必补，也毋须缝，一任自然，省掉烦恼，「向佛烧柱香儿，闲眠一觉」。何等潇洒，

\vfill
\begin{BookNotes}
\BookNote{42.}{蒋士铨，〈题杂家书画册子〉，引自王其和点校篡注，《板桥论画》（济南，山东画报出版社，2009年4月），页205。}
\end{BookNotes}

% 原书第 333 页（PDF 第 346 页）
\BookPage{333}{346}
\noindent 真是旷达。这首词虽是咏僧家，实则自况，人间俗事，推廓摒开，过着自适自在的生活。43〈道情〉有一首以「老头陀」为例，也表达同样的心境。曰：

\begin{BookQuote}
老头陀，古庙中，自烧香，自打钟，兔葵燕麦闲斋供。山门破落无关锁，斜日苍黄有乱松，秋星闪烁颓垣缝。黑漆漆蒲团打坐，夜烧茶炉火通红。〈道情·老
\end{BookQuote}

头陀》

尽管庙是古的，墙是破的，门是坏的，松是乱的，阳光也是斜的。兔葵燕麦，人神共食，生活相当清苦。但是老头陀自得其乐，自己生活自己安排，「日烧香，自打钟」，无须他人过问，优游闲适，怡然自得。虽然郑板桥欣赏和尙的生活，但是他不会出家为僧，功名、女人和酒，他还是留恋的。44由于郑板桥与佛教渊源深远，关系良善，因此，长期以来，都在维护佛法，帮佛教辩解。兹举一例，即可得悉。曰：

\begin{BookQuote}
佛自周昭王时下生，迄于灭度，足迹未尝履中国土。后八百年而有汉明帝说谎说梦，惹出这场事来，佛实不闻不晓。今不责明帝，而齐声骂佛，佛何辜乎？况自昌黎辟佛以来，孔道大明，佛焰渐息，帝王卿相，一遵六经四子之书，以为齐家治国平天下之道，此时而犹言辟佛，亦如同嚼蜡而已。和尚是佛之罪人，杀盗淫妄，贪婪势利，无复明心见性之规。秀才亦是孔子罪人，不仁不智，无礼无义，无复守先待后之意。秀才骂和尚，和尚亦骂秀才。语云：各人自扫门前雪，莫管他家屋瓦霜。〈焦山读书寄四弟墨书〉
\end{BookQuote}

上述文字，看出郑板桥明显的偏袒佛教，关心和尚。其一，佛教既已东来，国人要信佛，何必骂佛；其二，佛学既已昌盛，佛教渐息，无此道理；其三，佛门虽有作恶之僧徒，但儒家亦有无耻之人，彼此半斤八两；其四，儒佛二家各自研修相互尊重，不要批评攻讦。因此，郑板桥看似公平对待双方，但归本儒家的他，还是比较护持佛教。

这就是郑板桥的「无佛之佛」。

\vfill
\begin{BookNotes}
\BookNote{43.}{廖玉惠选注，《一竿烟雨一郑板桥诗词文选》，页221-222。}
\BookNote{44.}{丁家桐，《绝世风流一郑燮传》，页59-60。}
\end{BookNotes}

% 原书第 334 页（PDF 第 347 页）
\BookPage{334}{347}
四结语

郑板桥在官场中，崎岖难行，很不顺遂；在生活上，没落乍薄，辛苦清贫；在家居中，人丁不旺，屡遭变故；加上强烈执着的性格与忌恶如仇的特质，照说应很难跟人互动沟通、和乐善处，但他却打破这种隔阂与限制，反而开创一片天空，营造文人情境，为当时所推许，更享后世的仰望。细加探索，其与佛教的渊源与关系，应居要功。换言之，虽然郑板桥没有皈依三宝，但佛教与佛学却是他生命中不可或缺的滋养。

至于郑板桥之能与佛教建立殊胜因缘，可得有四：其一，郑板桥祖先埋葬佛寺坟地，且叔组福国出家为僧，由是播下佛缘种子。其二，郑板桥长期在佛寺读书、寄住、会友、卖画，使其得以解决生活诸多问题，进而得到心灵安慰。其三，郑板桥结交很多学佛的朋友，包括居士与僧侣，前者如扬州八怪及其他友人，至于僧侣可有多位，遍及各名山古刹。其四，僧友因清高孤绝，不落俗尘，因而与郑板桥极为投契，他们也乐于跟郑板桥交往，论佛说禅。

要之，郑板桥由于本诸儒家思想，接受民间习俗，因而没有托付佛陀，虽是信仰无佛，却时时有佛，处处见佛，无以名之，故曰「无佛之佛」。
